% !TeX program = pdflatex
\documentclass[11pt]{article}

% -----------------------
% Packages (barebones)
% -----------------------
\usepackage[margin=1in]{geometry}
\usepackage{setspace}
\usepackage{microtype}
\usepackage{amsmath, amssymb, amsthm, mathtools}
\usepackage{bm}              % bold math symbols
\usepackage{graphicx}
\usepackage{booktabs}
\usepackage{hyperref}
\usepackage[backend=biber]{biblatex}
\addbibresource{Proposal_bib.bib}

\onehalfspacing
\hypersetup{
  colorlinks=true,
  linkcolor=blue,
  citecolor=blue,
  urlcolor=blue
}

% -----------------------
% Metadata macros
% -----------------------
\newcommand{\Title}{Returns to Representation: Girls' Youth Sports Participation Following the Onset of the USWNT and WNBA}
\newcommand{\Author}{Tate Mason}
\newcommand{\Date}{\today}

% -----------------------
% Section helpers
% -----------------------
% Use: \Section{Motivation} ... \Section{Methods}
\newcommand{\Section}[1]{\section{#1}}
\newcommand{\Subsection}[1]{\subsection{#1}}
\newcommand{\Subsubsection}[1]{\subsubsection{#1}}

% -----------------------
% Equation helpers
% -----------------------
% Inline: \(\E[Y|X]\), display: \eqn{ ... }, numbered: \eqnlabel{key}{...}
\newcommand{\eqn}[1]{\begin{equation*}#1\end{equation*}}
\newcommand{\eqnlabel}[2]{\begin{equation}\label{eq:#1}#2\end{equation}}
\newcommand{\Eqref}[1]{(\ref{eq:#1})}

% Common operators
\DeclareMathOperator*{\argmax}{arg\,max}
\DeclareMathOperator*{\argmin}{arg\,min}

% Probability / expectation / variance
\newcommand{\Prob}{\mathbb{P}}
\newcommand{\E}{\mathbb{E}}
\newcommand{\Var}{\mathrm{Var}}
\newcommand{\Cov}{\mathrm{Cov}}

% Sets, reals, indicators
\newcommand{\R}{\mathbb{R}}
\newcommand{\1}[1]{\mathbb{I}\{#1\}}

% Vectors / matrices
\newcommand{\vct}[1]{\bm{#1}}     % bold vector
\newcommand{\mtx}[1]{\bm{#1}}     % bold matrix
\newcommand{\trans}{^\mathsf{T}}  % transpose
\newcommand{\inv}{^{-1}}

% Norms, inner products
\newcommand{\norm}[1]{\left\lVert #1 \right\rVert}
\newcommand{\ip}[2]{\left\langle #1, #2 \right\rangle}

% Derivatives
\newcommand{\dd}[2]{\frac{d #1}{d #2}}
\newcommand{\pp}[2]{\frac{\partial #1}{\partial #2}}
\newcommand{\ppn}[3]{\frac{\partial^{#3} #1}{\partial #2^{#3}}}

% Convenience delimiters
\newcommand{\paren}[1]{\left( #1 \right)}
\newcommand{\brak}[1]{\left[ #1 \right]}
\newcommand{\set}[1]{\left\{ #1 \right\}}

% Theorem-like envs (optional)
\theoremstyle{plain}
\newtheorem{assumption}{Assumption}
\newtheorem{proposition}{Proposition}

% -----------------------
% Title block
% -----------------------
\begin{document}

\begin{center}
  {\large \textbf{\Title}}\\[0.5em]
  {\large \Author}\\
  {\Affil}\\[0.25em]
  {\Date}
\end{center}

\vspace{0.5em}

% -----------------------
% Proposal structure
% -----------------------
\Section{Motivation and Research Question}
  In the late 1990's, young female athletes would be exposed to two new sources of professional athletics:
  the Women's National Basketball Association (WNBA) and the United States Women's National Team (USWNT).
  These entities provided young basketball and soccer players, respectively, the opportunity to envision
  a future in professional athletics. Having a pathway in mind, would these girls be spurred to increase
  participation in sports?

  My idea is to understand how the advent of WNBA and the popularization of the USWNT in the 1990's impacted 
  youth sports participation for girls. Beyond the preliminary question of the effect on youth sports, I am also
  interested in seeing the effect on college attendance, labor market outcomes, and health outcomes. College 
  attendance is of interest due to athletic scholarships being a potential byproduct of youth sports participation
  and labor market outcomes are of interest given those who attend college have better outcomes than those who do not.
  Finally, health outcomes are of interest since, even if a majority of these girls will not reach a professional level,
  there is evidence of long term health benefits associated with playing organized sports as a child \parencite{logan_youth_2020}.

\Section{Contribution}
  There is very little litereature on youth sports participation for girls. \textcite{stevenson_beyond_nodate}, covers a similar
  topic, but focuses instead on the implementation of Title IX on girls' participation. Her study finds a 10 percentage-point increase
  in sports participation, as well as a one percentage-point increase in female collge attendance and a one-to-two percentage-point increase
  in labor force participation.

  I would add to a study like this by focusing on a more contemporary event through professional sports' impact. Further, Stevenson does not
  spend time on health impacts. This would be an interesting avenue to add to a fresher look at educational attainment and labor market outcomes.

\Section{Methodological Strategy}
  The primary dimensions of variation for this topic are time and space. For time, there is staggered treatment, with the WNBA expanding from 8 teams
  in 1997 to 16 in 2000. For the USWNT, there are two sources of timing to exploit. First, in 1991, the USWNT won gold at the Women's World Cup in China,
  a hugely publicized accomplishment in the US. Secondly, the United States would host the 1999 Women's World Cup. This is insteresting, as both the effect
  of proximity and of representation can be tested.

  Geographically, the WNBA has set cities which host teams. This would, again, allow for testing of intensity of treatment by location. Similarly, for the WWC,
  having the US hosting presents the opportunity to examine the intensity of treatment in the eight host cities. Four of these cities happen to overlap. Those being
  Los Angeles, New York City, Portland, and an agglomerated area of DC-Baltimore. This further presents interesting opportunity for variation.

  I think doing a differences in differences in differences design in which I difference first having either of the sources of representation and having none. Then,
  differencing that from the instance of having both sources. A difficulty with this approach is the staggered treatment. For the WNBA, 2-4 teams were added yearly in
  the time horizon of interest. For the USWNT, there is the issue of two major sources of publicity across 8 years. Further, while I think endogeneity is not a huge
  threat here, as individuals did not have the choice of whether or not to live in a city which would be chosen to host the World Cup or a WNBA team, I worry about the
  proclivity to place a team in the largest cities or the cities whose demographics would best support a team/game. These are problems which will warrant more thought,
  and can hopefully be overcome.

\Section{Data}
  For data, I plan to follow the approach of \cite{stevenson_beyond_nodate}. The National Federation of State High School Associations (NFSH) keeps data on high school sports 
  participation by state and gender. This is not as granular as I'd hope, but is at least a starting point. IPUMS USA data can provide information on labor, location, and educational outcomes,
  as well as act as a panel with seeing individuals repeatedly. There is a gap here, however, in which I would have to calculate sports participation for a larger state, apply it to all
  women from IPUMS, and extrapolate. While this is my best option as of now, I am continuously looking for better options like the NSLY97 cohort, or some other data source to be found.

  Access for all data listed is public or already obtained. I have a lot of experience with IPUMS, and thus do not worry much about being able to use it effectively for labor and education
  outcomes. NFHS data I have not worked with, but it appears all I would need to do is scrape PDF participation surveys posted on their website. This should also be no problem, with the 
  biggest cost coming in the form of learning to scrape PDF's. 

  As stated above, I am still actively looking for fallbacks. I think data-wise, this project is mostly standard. The main problem is the disharmony between the datasets, but hopefully this
  can be remedied as I continue to look through data related to this field.
  

\printbibliography
  

\end{document}
